\chapter{Discussion}

The model selection results highlight the pivotal influence of spatial structure and environmental variables on \textit{Parthenium hysterophorus} distribution in Pakistan. Adding spatial smooth terms greatly improved model fit, evidenced by higher deviance explained and lower AIC, demonstrating that accounting for spatial autocorrelation is critical, as non-spatial models performed poorly by comparison.

While environmental predictors are important, they do not fully explain distribution patterns. Notably, increasing the spatial smooth basis from $k = 100$ to $k = 200$ substantially enhanced the model's ability to capture fine-scale spatial variation, suggesting that \textit{Parthenium}'s distribution is shaped by complex, geographically structured factors beyond those measured. Even after including key variables, altitude, climate, habitat, and road type, a significant residual spatial component persisted, as shown by the strong Gaussian process smooth. This suggests unmeasured or spatially structured drivers, such as local land use, dispersal dynamics, or microclimate, also play important roles and merit further study [X].

Altitude and road type were especially influential predictors, with model performance declining notably when either was excluded. The importance of road type is consistent with literature showing that transportation networks promote invasive species spread by offering dispersal corridors and disturbance regimes [X]. These insights underscore the need for targeted management along roadways and in vulnerable habitats [X].

The final REML-fitted GAM explained over 60\% of the deviance, confirming strong explanatory power. Yet, substantial residual spatial structure remains even after accounting for environmental and anthropogenic variables, indicating that other, finer-scale factors may influence distribution. Future research should incorporate higher-resolution spatial data, additional ecological covariates, and biotic interactions to clarify these drivers. Exploring \textit{Parthenium}'s use of allelochemicals to gain pollinator preference and enhance spread may also reveal important survival mechanisms [X].

The GAM results reveal that \textit{Parthenium hysterophorus} occurrence is primarily determined by the interplay of habitat disturbance, transport infrastructure, and climatic conditions [X]. The species is most likely to occur in open, disturbed habitats and along major road corridors, highlighting the strong association between invasion risk, land-use intensity, and human-mediated dispersal. Additionally, the notable nonlinear effects of climatic variables indicate that invasion is governed not just by disturbance, but by how these opportunities intersect with environmental suitability.

Habitat type is a major driver of invasion risk for \textit{P. hysterophorus}. Disturbed and unmanaged open habitats have the highest baseline occurrence, supporting the view of \textit{P. hysterophorus} as a classic disturbance-adapted pioneer species [X]. These plants quickly colonise areas where vegetation is disrupted, competition is low, and direct soil contact is possible. Following Grime's ecological framework, \textit{Parthenium} exhibits pioneer traits by prioritising rapid growth, flowering, and prolific seed production over long-term competition [X]. Agricultural Disturbed and Semi-natural Open habitats create especially favourable conditions with frequent disturbances, exposed soil, and resource pulses [X]. Such settings may also allow contaminants to affect soil composition and fertility, especially without biological fertilisers or other management interventions [X]. By contrast, habitats with active management or stronger natural resistance, such as Agricultural Managed land, Woody Vegetation, and Wet Habitats, exhibit much lower occurrence, indicating that both human control and biophysical barriers can suppress \textit{Parthenium} establishment.

Lower occurrence in managed and structurally complex habitats provides further ecological insight. In actively managed agricultural systems, repeated interventions, such as weed control, crop rotation, and maintaining managed ground, reduce establishment opportunities for \textit{Parthenium}. In woody vegetation, closed canopies limit light and established perennial species increase competition, collectively suppressing \textit{Parthenium} growth [X]. In wet habitats, excess moisture or related ecological factors may constrain invasion, either by directly causing physiological stress to \textit{Parthenium} or by enhancing competition from moisture-adapted plant communities [X]. These patterns illustrate how both management practices and natural habitat complexity serve as effective barriers to invasion.

Recent thermal accumulation, measured as 30-day growing degree days (GDD), was a strong but nonlinear predictor of \textit{Parthenium hysterophorus} occurrence. Rather than increasing steadily with temperature, occurrence peaked at intermediate to upper-intermediate GDD levels, particularly in less-managed habitats. This likely reflects the temperature dependence of multiple life stages, including germination, growth, flowering, and seed production, each with different optimal conditions [X]. Temperature therefore influences seasonal opportunities for establishment and spread, while excessive heat may cause sun bleaching, reducing photosynthesis and limiting growth and reproduction [X].

The bimodal or trimodal GDD pattern likely reflects seasonal rather than multiple discrete physiological optima. Occurrence peaked in August, with a secondary increase in November, while cooler months showed lower occurrence. These peaks may represent periods favourable for establishment and reproduction, with August corresponding to the main growth phase and November potentially reflecting secondary recruitment or post-monsoon persistence [X]. However, the apparent peaks should not be interpreted as sharp biological thresholds. Sparse observations and wide confidence intervals at thermal extremes increase uncertainty, and GAMs can produce exaggerated curvature where data are limited [X]. Interpretation should therefore focus on the better-sampled central portion of the gradient.

More broadly, GDD represents only one component of the \textit{P. hysterophorus} invasion niche. Moisture, disturbance, irrigation, hydrology, vegetation, and land management also influence occurrence and may covary with temperature [X]. Precipitation showed a nonlinear association, with occurrence highest under low-to-moderate rainfall and declining under wetter conditions. Thus, \textit{P. hysterophorus} appears most favoured where warmth coincides with intermediate moisture, particularly in open, disturbed habitats [X].

This distinction is also important for future projections. Under the current temperature-based model, \textit{Zygogramma bicolorata} remains broadly within climatic tolerance across India, while \textit{Parthenium} suitability contracts by 2050. However, the model excludes hydrological buffering. Irrigation studies in Pakistan show that incorporating water availability can substantially increase predicted \textit{Parthenium} suitability, particularly in marginal areas [X]. Temperature-only projections may therefore overestimate future declines where moisture offsets thermal stress [X].

Model forecasts indicate that, under the current assumptions, the beetle largely remains within climatic tolerance across India, while suitable climate for \textit{Parthenium hysterophorus} contracts over time. This apparent beetle resilience reflects the permissive annual thermal proxy used. By 2050, reduced overlap is primarily driven by declining \textit{Parthenium} suitability as temperatures exceed physiological thresholds that may impair metabolic processes and persistence [X]. These projections should therefore be interpreted as baseline thermal risk assessments rather than comprehensive ecological forecasts.

However, the temperature-only model omits key environmental drivers, particularly precipitation and moisture, which can buffer heat stress and support persistence under otherwise unsuitable temperatures. Irrigation studies in Pakistan show that incorporating water availability can substantially increase predicted \textit{Parthenium} suitability, particularly in marginal areas [X]. Temperature-only projections may therefore overestimate future declines where hydrological buffering occurs [X].

The 2050 heat-stress panel provides an important seasonal perspective, revealing thermal bottlenecks obscured by annual averages. Although annual conditions remain broadly permissive for \textit{Zygogramma bicolorata}, large areas of central and southern India may experience summer temperatures associated with heat stress or thermal unsuitability, potentially limiting beetle persistence before the monsoon [X]. Cooler regions, particularly the northeast and areas below 38$^{\circ}$C, may provide thermal refugia. However, these projections use coarse macroclimatic data and do not capture cooler microclimates provided by shaded foliage or moist soils [X].

The heat-stress assessment also excludes monsoon timing, soil moisture, diapause, and seasonal emergence dynamics, all of which influence \textit{Z. bicolorata} persistence. In addition, real-world biocontrol involves dynamic ecological feedbacks: beetle populations may increase following \textit{Parthenium} expansion but subsequently decline as food resources are depleted, while predators, parasitoids, and phenological mismatches may further constrain control [X] [X]. Accordingly, the figure provides a broad macroclimatic assessment rather than a process-based biocontrol forecast. Its main contribution is to identify temperature as an important constraint, highlight potential future narrowing of weed--beetle overlap, and demonstrate that summer heat extremes may create bottlenecks that annual climate averages overlook [X].

Comparison of the Pakistan-calibrated GAM and India-calibrated MaxEnt models revealed only partial spatial agreement across Pakistan. Weak positive correlations (Spearman's $\rho$; Pearson's $r$) indicate that the models differ in both rank ordering and predicted suitability, despite being compared across the same spatial support.

This divergence likely reflects differences in calibration domain, predictors, model framework, and data. The GAM was calibrated using Pakistani data, whereas MaxEnt was trained on Indian data and transferred to Pakistan, meaning the models represent different environmental and realised-niche relationships. MaxEnt uses broad climatic predictors, while the GAM incorporates growing degree days, precipitation, altitude, and habitat cover. The GAM also restricted projections to its training environmental envelope and excluded spatial smoothing to minimise unsupported extrapolation. Differences between controlled presence--absence data for Pakistan and presence-only data with varying sampling bias for India may further contribute to the contrasting suitability surfaces.

Despite weak overall concordance, both models identified a shared southern hotspot. Areas classified as ``both high'' therefore provide more robust candidate regions because they are supported across two independent calibration contexts. In contrast, ``GAM only high'' and ``MaxEnt only high'' areas represent model uncertainty and may reflect differences in predictor sensitivity, local environmental structure, or extrapolation [X].

The low Jaccard overlap among the top 25\% of suitability cells reinforces that hotspot identification is model-dependent. However, comparison on an identical aligned valid-cell set and use of relative suitability ranks improve comparability between models. The rank-difference surface further identifies where each model assigns greater suitability [X].

Overall, Figure~\ref{fig:pakistan-concordance} demonstrates that transferred MaxEnt predictions capture some, but not all, of the spatial pattern identified by the Pakistan-calibrated GAM. Areas of agreement provide useful targets for future validation, while divergent regions highlight uncertainty and the limitations of transferring models across calibration geographies [X].

\section{Conclusion}

Overall, the GAM captured fine-scale, season- and habitat-linked occurrence risk, while MaxEnt identified broader climatic gradients. Both models highlighted disturbed, warmer, invasion-prone areas, but differed in hotspot locations and extent, underscoring that \textit{Parthenium} favours specific environments requiring further investigation. For management, derelict lands, semi-natural open habitats, and disturbed agricultural areas should be prioritised for surveillance and control, with interventions timed to peak seasonal windows and guided by real-time climate data [X]. Public awareness remains limited: hand-pulling is common in small fields, but knowledge of biocontrol is low [X]. Broader dissemination of biocontrol information and stronger prevention policies, such as stricter border and transport surveillance, are needed [X]. Given that conflict and increased transport may accelerate seed dispersal across South Asia, management must address both anthropogenic pathways and climatic suitability. Species distribution models remain scenario-based approximations, limited by predictor choice and unable to fully resolve dispersal, biotic interactions, microclimate, or feedback [X]. Nevertheless, this study provides a crucial spatial baseline for assessing climatic risk, identifying management priorities, and guiding research on how warming may affect \textit{Parthenium} and its biocontrol agent in Pakistan [X].